\chapter{Основные понятия}
\label{ch:concepts}

Это самая важная глава для мастера. Платформа различает четыре вещи, которые в разговоре обычно называют одним словом ``игра'': сценарий, запущенный мир, подход и кампанию. Интерфейс эту разницу почти не подсказывает, поэтому разберём её здесь.

\section{Цепочка от замысла до игры}
\label{sec:concepts:chain}

\begin{center}
\texttt{Сценарий} $\rightarrow$ \texttt{Запущенный мир} $\rightarrow$ \texttt{Лобби} $\rightarrow$ \texttt{Подход} $\rightarrow$ (\texttt{Кампания})
\end{center}

\paragraph{Сценарий}
Заготовка: локации, персонажи мастера, предметы, сюжетные узлы. Живёт на \texttt{/scenarios}. Это \emph{шаблон} --- играть прямо в нём нельзя, и правки в нём не влияют на уже идущие игры.

\paragraph{Запущенный мир}
Сценарий, переведённый в игровое состояние (``launched''). По сути --- слепок заготовки, который дальше живёт своей жизнью: персонажи в нём получают урон, счётчики крутятся, предметы переходят из рук в руки. Список --- на \texttt{/launched-scenarios}.

Ключевое следствие: один сценарий можно запустить несколько раз, для разных компаний игроков, и эти миры не будут знать друг о друге.

\paragraph{Лобби}
Комната ожидания перед игрой. Игроки подключаются, мастер раздаёт им персонажей и запускает игру. Адрес --- \texttt{/lobby/<id>}.

\paragraph{Подход}
Одна игровая встреча внутри запущенного мира --- то, что за столом называют ``сессия''. Открывается по адресу \texttt{/session/<id>}. Мир может пережить много подходов подряд.

\paragraph{Кампания}
Объединяет подходы в связную историю. Главное, что она даёт, --- \emph{перенос состояния} (carryover): когда подход завершается, состояние персонажей, инвентарь и счётчики переезжают в следующий. Без кампании каждый подход начинается с исходного состояния мира. Раздел --- \texttt{/campaigns}.

\screenshotOptional{user/campaigns.png}{раздел «Кампании»}

Сам сайт описывает кампании так: «многошаговые истории с общим миром: один запущенный сценарий живёт между подходами и партиями, персонажи и мир переносятся между эпизодами». Если это описание вам понятно --- главу можно считать усвоенной.

\section{Сущности сценария}
\label{sec:concepts:entities}

Наполнение сценария состоит из объектов нескольких типов. Подробный разбор каждого --- в главе \ref{ch:objects}, здесь коротко:

\begin{description}
\item[Локация] Место действия. Может нести карту-картинку с областями, вложенные локации.
\item[Story beat] Сюжетный узел: что происходит, кто участвует, что игроки услышат.
\item[NPC] Персонаж под управлением мастера.
\item[Предмет] Вещь, которую можно найти, отдать, использовать.
\item[Персонаж игрока] Заготовка героя. Во время игры привязывается к конкретному человеку.
\item[Заметка] Текст для мастера или для игроков.
\item[Счётчик] Трекер: часы, ресурсы, шкала угрозы --- что угодно с числом.
\item[Препятствие] Проверка: метод, сложность, что при успехе и при провале.
\item[Задача] Пункт в списке подготовки сценария, для самого мастера.
\end{description}

\paragraph{Формы разные у разных систем}
Набор полей у персонажа в Dungeon World и в Blades in the Dark отличается, потому что форму задаёт выбранная система правил, а не сайт. Поэтому не удивляйтесь, что в одном сценарии у персонажа есть ``характеристики'', а в другом --- нет.

\section{Что игроки видят, а что нет}
\label{sec:concepts:exposure}

Разделение на ``знает мастер'' и ``знают игроки'' проходит через всю платформу.

\paragraph{Заготовка сцены}
У локаций и story beat есть описание того, что игроки увидят, когда мастер откроет эту сцену. Оно пишется заранее, при подготовке сценария, и не показывается раньше времени.

\paragraph{Туман войны на карте}
Карта локации размечается областями. Мастер открывает их по одной --- игроки видят только раскрытые.

\paragraph{Показ объекта}
Мастер может ``показать'' конкретный объект --- NPC, предмет, локацию --- и он крупно выведется у всех участников. Удобно, когда в сцену входит новое лицо.

\paragraph{Адресные сообщения}
Помимо общего чата есть личные рассылки конкретному игроку --- то, что раньше передавали запиской через стол. Игрок может на них ответить.

\section{Системы правил}
\label{sec:concepts:rulesystems}

Система правил выбирается при создании сценария и потом не меняется. От неё зависят поля сущностей и то, как работают броски.

\begin{center}
\begin{tabular}{ll}
\textbf{Система} & \textbf{Насколько готова} \\
\hline
Dungeon World & Проработана полностью \\
PbtA (базовый слой) & Работает, основа для остальных \\
Gumshoe (и три сеттинга) & Сущности есть, механика действий слабее \\
Blades in the Dark & Каркас \\
Everyone is John & Каркас \\
I'm Bitter and I Keep Records & Каркас \\
\end{tabular}
\end{center}

\paragraph{Практический совет}
Если хотите, чтобы работало всё, включая ходы, урон и автоматический подсчёт --- берите Dungeon World. На остальных системах вести сценарий можно, но броски и последствия придётся считать за столом самим.
